%! Special Directions --- Setting Up Eigenvalues — Unit 6, Lesson 6.0 — Warm-Up (Answer Key)
\documentclass[10pt]{article}
\usepackage{linalg-article}
\usepackage{linalg-key}   % pulls in -boxes; do NOT also load -boxes
\newcommand{\TallMath}[1]{$\displaystyle #1\rule[-1.4em]{0pt}{3.2em}$}

\begin{document}

\pageheader{Unit 6, Lesson 6.0}{Warm-Up --- Answer Key}
\namedateperiod

\vspace{0.12in}

\begin{notesbox}{Getting Ready: Ideas We'll Use Today}
\small These three ideas from Units~1 and~5 set up today's ``special directions.'' Answer quickly.

\vspace{4pt}
\textbf{1. Matrix times a vector (Unit 1).} Compute $A\mathbf{x}$ for
$A=\begin{bmatrix} 2 & 1 \\ 1 & 2 \end{bmatrix}$ and $\mathbf{x}=\begin{bmatrix} 1 \\ 1 \end{bmatrix}$:
\[
  A\mathbf{x}=\begin{bmatrix} 2 & 1 \\ 1 & 2 \end{bmatrix}\begin{bmatrix} 1 \\ 1 \end{bmatrix}
  = \begin{bmatrix}{\color{keyred}\mathbf{3}}\\[2pt]{\color{keyred}\mathbf{3}}\end{bmatrix}.
\]

\vspace{4pt}
\textbf{2. Scalar multiples (Unit 1).} Is the vector $\begin{bmatrix} 3 \\ 3 \end{bmatrix}$ a
\emph{multiple} of $\begin{bmatrix} 1 \\ 1 \end{bmatrix}$? If so, what number do you multiply by?
\[
  \begin{bmatrix} 3 \\ 3 \end{bmatrix} = {\color{keyred}\mathbf{3}}\cdot\begin{bmatrix} 1 \\ 1 \end{bmatrix}.
\]

\vspace{4pt}
\textbf{3. The $2\times2$ determinant (Unit 5).} Recall $\det\begin{bsmallmatrix} a & b \\ c & d
\end{bsmallmatrix}=ad-bc$. Compute
\[
  \det\begin{bmatrix} 2 & 1 \\ 1 & 2 \end{bmatrix} = {\color{keyred}\mathbf{(2)(2)-(1)(1)=3}}.
\]
\end{notesbox}

\begin{teachernote}
All three items set up eigenvectors. Items~1 and~2 \emph{are} the lesson in miniature: $A$ times
$\begin{bsmallmatrix}1\\1\end{bsmallmatrix}$ gives $\begin{bsmallmatrix}3\\3\end{bsmallmatrix}=
3\begin{bsmallmatrix}1\\1\end{bsmallmatrix}$, so the matrix did nothing to that direction but
\emph{multiply it by $3$} --- $\begin{bsmallmatrix}1\\1\end{bsmallmatrix}$ is an eigenvector and $3$
is its eigenvalue. Item~3 rehearses the $2\times2$ determinant, the tool that will \emph{find} the
eigenvalues in Lesson~6.1. Debrief items~1--2 aloud before starting the notes.
\end{teachernote}

\end{document}
